\chapter{Results, Evaluation, and Security Analysis}

\section{Experimental Architecture}
To rigorously validate the computational efficacy and operational security guarantees of the condition-wise blockchain healthcare design, automated evaluation scripts were engineered under the \texttt{experiments/} folder. In Phase 2.2, the experiments were extended to support two input modes: \texttt{single} (a backward compatible single-text workload) and \texttt{patient\_docs} (a dataset-style workload generated by \texttt{patient\_data.py}). For \texttt{patient\_docs}, record identifiers were standardized using the persistent disease-to-id mapping in \texttt{data/disease\_code\_map.json}, producing record keys in the format \texttt{<patient\_number>\_<disease\_id>}. All results are exported as CSV files under \texttt{runtime\_experiments/}.

\section{Latency Breakdown Analysis}
Computational overhead generated by distributed computing boundaries is heavily scrutinized in blockchain systems. Evaluating the system iteratively provided insightful metrics breaking down the cumulative operation into independent timeline vectors representing network queries, node unwrap mechanisms, mathematical operations, and decryption efforts.

\begin{table}[ht]
\centering
\begin{tabular}{|l|c|c|c|c|c|c|}
\hline
\textbf{Priority Level} & \textbf{Threshold ($k$)} & \textbf{Fabric (ms)} & \textbf{Unwrap (ms)} & \textbf{Reconstruct (ms)} & \textbf{Store Get (ms)} & \textbf{Total (s)} \\
\hline
HIGH & 2 & 42.13 & 16.19 & 0.20 & 9.28 & \textbf{2.160} \\
MEDIUM & 3 & 27.22 & 26.05 & 0.45 & 9.64 & \textbf{2.159} \\
LOW & 4 & 17.41 & 39.00 & 0.55 & 13.78 & \textbf{2.154} \\
\hline
\end{tabular}
\caption{Representative Latency Breakdown per Data Priority (Phase 2.2, \texttt{patient\_docs} mode)}
\label{tab:latency}
\end{table}

As evident in Table \ref{tab:latency}, the dominant costs arise from peer coordination and off-chain retrieval, while the core cryptographic computations (Shamir reconstruction and AES-GCM decryption) remain small relative to the end-to-end reconstruction time. Importantly, priority selection directly controls the threshold $k$, which affects how many peers must participate during reconstruction.

\section{Fault Tolerance Efficacy}
Robust network operability amidst catastrophic server failure profiles outlines the primary justification for integrating threshold cryptography techniques initially. We conducted simulated network failure tests—selectively killing peer processes in distinct combinations—to confirm Shamir bounds. System resiliency accurately reflected expectations across diverse $k:n$ structures out of $N=5$ participants. 

\begin{table}[ht]
\centering
\begin{tabular}{|l|c|c|c|c|}
\hline
\textbf{Data Priority} & \textbf{Threshold ($k$)} & \textbf{Failed Nodes} & \textbf{Active Nodes} & \textbf{Success Ratio} \\
\hline
HIGH & 2 & 0 & 5 & 100\% \\
HIGH & 2 & 1 & 4 & 100\% \\
HIGH & 2 & 2 & 3 & 100\% \\
HIGH & 2 & 3 & 2 & 100\% \\
HIGH & 2 & 4 & 1 & \textbf{0.0\%} \\
\hline
MEDIUM & 3 & 0 & 5 & 100\% \\
MEDIUM & 3 & 1 & 4 & 100\% \\
MEDIUM & 3 & 2 & 3 & 100\% \\
MEDIUM & 3 & 3 & 2 & \textbf{0.0\%} \\
\hline
LOW & 4 & 0 & 5 & 100\% \\
LOW & 4 & 1 & 4 & 100\% \\
LOW & 4 & 2 & 3 & \textbf{0.0\%} \\
\hline
\end{tabular}
\caption{System Capability against Incremental Peer Node Failures (Phase 2.2, $N=5$)}
\label{tab:fault_tolerance}
\end{table}

Data aggregated in Table \ref{tab:fault_tolerance} explicitly supports the Shamir threshold definition. High priority records are configured with a lower $k$, enabling reconstruction under more peer failures, while low priority records enforce stricter consensus (higher $k$), improving compromise resistance at the expense of fault tolerance.

\section{Compromise Resistance Analysis}
To complement fault tolerance, we also evaluated compromise resistance analytically by varying the number of compromised peers ($c$) and checking whether an attacker could accumulate at least $k$ usable shares. For a $k$-of-$n$ threshold, the attacker can reconstruct the secret once $c \geq k$. Therefore, higher thresholds (used for lower priority records) increase resistance against partial peer compromise.

\begin{table}[ht]
\centering
\begin{tabular}{|l|c|c|c|}
\hline
\textbf{Priority Level} & \textbf{Threshold ($k$)} & \textbf{Min. Compromised Peers for Attack Success} & \textbf{Attacker Success Condition} \\
\hline
HIGH & 2 & 2 & $c \geq 2$ \\
MEDIUM & 3 & 3 & $c \geq 3$ \\
LOW & 4 & 4 & $c \geq 4$ \\
\hline
\end{tabular}
\caption{Compromise Resistance Summary for $N=5$ Peers (Phase 2.2)}
\label{tab:compromise}
\end{table}

\section{Security Paradigms against Orderer Exploitation}
The strategic external decoupling of the Trusted Authority fundamentally thwarted hypothetical "Orderer-as-TA" security failure frameworks. Specifically, our implemented architecture mitigated extensive blast radiuses:
\begin{itemize}
    \item \textbf{Trust Boundary Separation:} A malicious insider compromising the core Fabric Ordering nodes only gained capabilities extending to data ordering and consensus disruption. Since cryptographic key wrapping protocols existed isolated off-chain across NMK logic boundaries, the encrypted Object database remained inherently invulnerable. 
    \item \textbf{Insider Threats:} Separation of duties strictly enforced operational governance. System administrators operating node topologies possessed completely disconnected permissions and token-capabilities compared to external hospital and healthcare administrative actors. 
    \item \textbf{Versioning Cryptanalysis Check:} Evaluation demonstrated that subsequent record updates utilized entirely randomized Nonce configurations mathematically proving that unauthorized possession of previous state ciphertexts effectively wielded zero potential for retro-computing future block states \cite{samy2020fog}.
\end{itemize}

Ultimately, exhaustive security parameters coupled alongside remarkable real-time functional metrics validated the profound capabilities of decentralized healthcare frameworks enhanced logically via condition-sensitive blockchain integrations \cite{shorfuzzaman2021blockchain}.
